% Layout copied from the source paper by scripts/extract_tables.py; numbers are \nv macros.
\begin{table}[H]
\centering\small
\caption{Second-round accounts and their prospective tests on held-out material.}
\label{tab:d4-r2tests}
\begin{tabularx}{\linewidth}{@{}l Y c c@{}}
\toprule
Tests & Account & Passed & Verdict \\
\midrule
R1--R6 & \qsix{}: \hd{0}{2} $\to$ \mlp{0} $\to$ \mlp{24}, which decides teen; \mlp{18} writes the last digit's suffix (4 and 6 digits) & \nv{d4.r2.path.pass} of \nv{d4.r2.path.n} & \nv{d4.r2.path.verdict} \\
S1--S3 & \qsix{}: later previous-token heads back up the copy, mostly at two digits & \nv{d4.r2.backup.pass} of \nv{d4.r2.backup.n} & \nv{d4.r2.backup.verdict} \\
O3--O8 & \olmo{}: a weak route across the token boundary, not a missing one (3--10 digits) & \nv{d4.r2.olmo.pass} of \nv{d4.r2.olmo.n} & \nv{d4.r2.olmo.verdict} \\
X & three more one-token-per-digit models: early previous-token heads carry the teen check & \nv{d4.r2.other.pass} of \nv{d4.r2.other.n} & \nv{d4.r2.other.verdict} \\
\bottomrule
\end{tabularx}
\end{table}
